\documentclass[12pt,a4paper]{article}
\usepackage[utf8]{inputenc}
\usepackage[T1]{fontenc}
\usepackage{amsmath}
\usepackage[left=2.00cm, right=2.00cm, top=2.00cm, bottom=2.00cm]{geometry}
\usepackage{amssymb}
\usepackage{dsfont}
\usepackage{graphicx}
\usepackage{ragged2e}
\usepackage{yhmath}
\usepackage[italian]{babel}
% title
\title{Foglio di esercizi\\Complementi di matematica per le scienze chimiche\\ Anno accademico 2025/26}
\author{Massimiliano Ghiotto}
\date{}
% comandi per quesito e soluzione
\usepackage{amsthm}
\theoremstyle{definition}
\newtheorem{sol}{Soluzione all'esercizio}
\newtheorem{exe}{Esercizio}

% insiemi numerici
\newcommand{\numberset}{\mathbb}
\newcommand{\N}{\numberset{N}}
\newcommand{\Z}{\numberset{Z}}
\newcommand{\Q}{\numberset{Q}}
\newcommand{\R}{\numberset{R}}
% cartella figure
\graphicspath{{figures/}}
\begin{document}
\maketitle
\noindent
Per qualsiasi domanda potete contattarmi per mail all'indirizzo:
\\
\textit{massimiliano.ghiotto01@universitadipavia.it}

\section{Statistica descrittiva}

\begin{exe}
	La tabella seguente rappresenta, al variare dell'anno, il numero di milioni di barili di olio che gli Stati Uniti hanno importato dalla Norvegia
	\begin{center}
		\begin{tabular}{|c|c|c|c|c|c|c|c|c|c|c|c|}
			\hline
			anno              & 1999 & 2000 & 2001 & 2002 & 2003 & 2004 & 2005 & 2006 & 2007 & 2008 & 2009 \\
			\hline
			milioni di barili & 96   & 110  & 103  & 127  & 60   & 54   & 43   & 36   & 20   & 11   & 22
			\\
			\hline
		\end{tabular}
	\end{center}
	Calcolarne:
	\begin{enumerate}
		\item  la media e la mediana;
		\item la varianza e la deviazione standard.
		\item calcolare i quartili e la distanza interquartile.
	\end{enumerate}
\end{exe}

\begin{exe}
	A $ 1000 $ (mille) studenti iscritti all'Università di Pavia viene proposto un sondaggio dove in particolare si chiede di esprimere un giudizio su un servizio offerto agli studenti, usando una scala da
	$ 1 $ a $ 6 $ ($ 1 $ = pessimo, $ 6 $ = eccellente). Le risposte ottenute sono riassunte qui sotto, dove per ognuno dei
	giudizi fra parentesi vi è il numero di voti ottenuti.
	\[
		1\ (187)\quad 2\ (82)\quad 3\ (260)\quad 4\ (98)\quad 5\ (159)\quad 6\ (214) \]
	\begin{enumerate}
		\item  Calcolare media, moda, mediana dei risultati;
		\item  determinare anche varianza e deviazione standard;
		\item  fare un istogramma per visualizzare meglio i dati.
	\end{enumerate}
\end{exe}

\section{Calcolo combinatorio}

\begin{exe}
	Cinque amici Ada, Beatrice, Carlo, Dario ed Emanuele si sfidano in una corsa di $ 1000 \ m$.
	\begin{enumerate}
		\item In quanti modi possono classificarsi i $ 5 $ amici?
		\item Quanti modi ci sarebbero sapendo che Ada è arrivata seconda?
		\item Quanti modi ci sarebbero se si sapesse soltanto che Ada è arrivata prima o seconda?
	\end{enumerate}
\end{exe}

\begin{exe}
	Nel gioco del lotto sono presenti i numeri da $ 1 $ a $ 90 $ (cioè sono presenti $ 90 $ numeri in totale), quartina e cinquina sono rispettivamente sequenze di  $ 4 $ e $ 5 $ numeri distinti.
	\begin{enumerate}
		\item Calcolare quante cinquine sono possibili.
		\item Calcolare la probabilità di fare cinquina (giocando cinque numeri).
		\item Calcolare la probabilità di fare quartina (giocando cinque numeri).
	\end{enumerate}
\end{exe}

\begin{exe}
	Il professor Layton si dimentica il codice di tre cifre del suo lucchetto per la bicicletta, decide allora di provare tutte le possibilità finché non trova quella corretta.
	\begin{enumerate}
		\item Calcolare la probabilità che indovini la combinazione del lucchetto al primo tentativo.
		\item Calcolare la probabilità che indovini la combinazione del lucchetto al primo o al secondo tentativo.
		\item Calcolare quante possibilità sapendo che nella combinazione che apre il lucchetto è presente almeno una cifra $ 3 $.
	\end{enumerate}
\end{exe}

\begin{exe}
	La signora Erina Joestar vuole ridipingere le $ 5 $ pareti ($ 4 $ pareti laterali e il soffitto) della sua camera ed ha a disposizione $ 6 $ colori: giallo, verde, bianco, azzurro, rosa, arancione. Sapendo che ogni parete potrebbe essere colorata con un colore distinto calcolare:
	\begin{enumerate}
		\item quanti sono i possibili modi di colorare la stanza?
		\item quanti sono i possibili modi per colorare la stanza sapendo che il soffitto ha deciso di farlo bianco?
		\item quanti sono i possibili modi per colorare la stanza sapendo che Erina ha deciso che almeno una parete sarà bianca.
	\end{enumerate}
\end{exe}

\begin{exe}
	Quante diagonali ha un poligono di $ n $ lati?
\end{exe}

\begin{exe}
	Un gruppo di $ 8 $ ragazzi vanno a vedere al cinema il nuovo film di Super Mario Bros, ai ragazzi viene assegnata un'intera fila del cinema da $ 15 $ posti. In quanti modi si possono disporre a sedere i ragazzi lungo la fila del cinema?
\end{exe}

\begin{exe}
	In una scacchiera $ 7 \times 7 $ si trova una pedina nel vertice in basso a sinistra. Le mosse consentita dalla pedina sono:
	\begin{itemize}
		\item spostarsi di una casella verso destra;
		\item  spostarsi di una casella verso l'alto.
	\end{itemize}
	Quanti percorsi possibili può percorrere la pedina per raggiungere il vertice in alto a destra della scacchiera?
\end{exe}

\begin{exe}
	In quanti modi si possono distribuire $ 10 $ caramelle a $ 3 $ bambini? Ed in quanti modo si possono distribuire se ogni bambino deve ricevere almeno $ 2 $ caramelle?
\end{exe}

\begin{exe}[$\star$]  Sia $ f:A \to B $ una funzione tra due insiemi $ A $ e $ B $.
	\begin{enumerate}
		\item Calcolare quante sono tutte le possibili \textbf{funzioni} sapendo che l'insieme $ A $ ha $ 5 $ elementi e l'insieme $ B $ ha $ 4 $ elementi.
		\item  Con gli stessi insiemi di prima calcolare quante sono le possibili \textbf{funzioni suriettive}.
	\end{enumerate}
\end{exe}

\begin{exe}
	Sia $ f: \mathbb{R}^n \to \mathbb{R} $ funzione di classe $ C^k $.
	\begin{itemize}
		\item Quante derivate di ordine esattamente pari a $ k $ si possono calcolare?
		\item Ricordare il teorema di Schwarz, quante sono le derivate parziali miste che coincidono con la derivata $ \frac{\partial f}{\partial x_1^3\ \partial x_2^2\ \partial x_4^4} $?
	\end{itemize}
\end{exe}

\section{Probabilità}

\begin{exe}
	Qual è la probabilità che lanciando due dadi escano:
	\begin{itemize}
		\item due $ 4 $;
		\item un $ 3 $ e un $ 5 $;
		\item due numeri pari;
		\item due numeri la cui somma è $ 9 $;
		\item due numeri uguali.
	\end{itemize}
\end{exe}

\begin{exe}
	Da un mazzo da $ 52 $ carte se ne estrae una. Calcolare la probabilità che sia:
	\begin{itemize}
		\item una carta di picche o una figura di cuori.
		\item una figura o una carta rossa.
	\end{itemize}
\end{exe}

\begin{exe}
	Una macchina produce pezzi meccanici e, su una produzione di $ 400 $ pezzi, $ 20 $ sono difettosi per peso, $ 30 $ per lunghezza e $ 360 $ sono perfetti. Calcola la probabilità che, prendendo a caso un pezzo:
	\begin{enumerate}
		\item sia difettoso;
		\item abbia entrambi i difetti;
		\item gli eventi $ E_1 = pezzo\ difettoso\ per\ peso $ e $ E_2 = pezzo\ difettoso\ per\ lunghezza $ sono indipendenti?
	\end{enumerate}
\end{exe}

\begin{exe}
	Un'urna contiene $ 8 $  palline rosse e $ 4 $ gialle. Calcola la probabilità che, estraendo consecutivamente due palline senza rimettere la pallina estratta nell'urna esse siano:
	\begin{enumerate}
		\item due palline rosse;
		\item due palline gialle;
		\item la prima rossa e la seconda gialla;
		\item una pallina rossa e l'altra gialla.
	\end{enumerate}
\end{exe}

\begin{exe}
	Vi è un'urna contenente $ 8 $ palline rosse, $ 5 $ blu, $ 7 $ bianche. Estraggo due volte,
	re-inserendo nell'urna la prima pallina estratta. Qual è la probabilità che alla
	prima estrazione esca una pallina rossa o bianca e alla seconda estrazione esca
	una pallina blu.
\end{exe}

\begin{exe}
	Ho un mazzo di $ 52 $ carte (da poker) ed estraggo direttamente $ 2 $ carte. Qual
	è la probabilità che le carte estratte siano:
	\begin{enumerate}
		\item due figure?
		\item una figura la prima e un asso la seconda?
	\end{enumerate}
\end{exe}

\begin{exe}
	Un commesso ha mescolato i barattoli di legumi sullo scaffale di un supermercato. Sappiamo che ci sono $ 7 $ barattoli di piselli, $ 9 $ barattoli di fagioli e $ 6 $ barattoli di lenticchie. Si prendono consecutivamente $ 3 $ barattoli. Calcola la probabilità che:
	\begin{enumerate}
		\item siano tutti di piselli;
		\item uno sia di piselli e due di fagioli;
		\item ce ne sia uno per tipo;
		\item non ci sia alcun barattolo di lenticchie.
	\end{enumerate}
\end{exe}

\begin{exe}
	Si hanno due urne. La prima contiene $ 4 $ palline bianche e $ 6 $ rosse, la seconda ne contiene $ 3 $ bianche e $ 5 $ rosse. Calcola la probabilità che, estraendo una pallina da ciascuna urna, esse siano:
	\begin{enumerate}
		\item entrambe bianche;
		\item bianca dalla prima urna e rossa dalla seconda urna;
		\item una bianca e una rossa.
	\end{enumerate}
\end{exe}

\begin{exe}[Paradosso dei compleanni]
	Calcolare la probabilità che in una classe di $ 40 $ persone almeno $ 2 $ siano nate nello stesso giorno.
\end{exe}

\section{Processo di Bernoulli}

\begin{exe}
	Un esame universitario consiste in un quiz a crocette con $ 10 $ domande a risposta multipla con $ 4 $ possibilità di risposta per ciascuna domanda e la risposta corretta è una sola. Si passa l'esame se in totale si compiono al massimo due errori. Qual è la probabilità di passare l'esame se si risponde a caso a tutte le domande?
\end{exe}

\begin{exe}
	Per un corso di studi universitario con laurea triennale, la probabilità che uno studente
	che si iscrive al primo anno si laurei effettivamente è $ p = 0.6 = \frac35 $. Qual è la probabilità che, su $ 10 $ studenti immatricolati,
	\begin{enumerate}
		\item  nessuno si laurei;
		\item uno si laurei;
		\item almeno uno si laurei;
		\item tutti si laureino.
	\end{enumerate}
\end{exe}

\begin{exe}
	Ho un mazzo di $ 40 $ carte (da briscola) ed estraggo per $ 5 $ volte una carta rimettendo ogni
	volta la carta estratta nel mazzo. Introdotta la variabile aleatoria “X = numero delle volte in cui esce un Re” e richiamato il concetto di distribuzione binomiale, calcolare
	\begin{enumerate}
		\item la probabilità $ P({X \le 1}) $, cioè che un Re esca al più una volta nelle estrazioni;
		\item la probabilità che un Re esca almeno una volta;
		\item la media (o valore atteso), la varianza e la deviazione standard (o scarto quadratico medio) per la variabile aleatoria $ X $.
	\end{enumerate}
\end{exe}

\begin{exe}
	Calcolare in modo esplicito la densità discreta della legge binomiale $ Bin(5, 0.15) $,
\end{exe}

\begin{exe}
	Se lancio $ 2 $ dadi la probabilità di fare $ 12 $ é $ \frac{1}{36} $: dunque mi aspetto che in $ 36 $ lanci esca un $ 12 $. Qual é la probabilità che ciò accada? Qual é il minimo numero di lanci da fare affinché la probabilità di uscita di un $ 12 $ sia almeno $ 0.5 $?
\end{exe}

\section{Processo di Poisson}

\begin{exe} Una folla di persone aspetta in fila indiana il taxi, alla stazione centrale. Sia $ X $ il numero di taxi che passa in un'ora e $ N $ il numero medio di texi che passano in un'ora ($ N $ è noto).
	\begin{enumerate}
		\item Sotto quali ipotesi sul fenomeno si può assumere che $ X $ sia una variabile aleatoria di Poisson?
		\item Ci si mette d'ora in poi nelle ipotesi del punto precedente. Se $ N = 20 $, qual è la probabilità che nella prossima mezz'ora passino esattamente $ 5 $ taxi?
		\item Nelle ipotesi precedenti qual è la probabilità che nella prossima mezz'ora passino non più di $ 5 $ taxi?
		\item  Qual'è la probabilità che nella prossima mezz'ora passino almeno $ 7 $ taxi?
	\end{enumerate}
\end{exe}

\begin{exe}
	Un apparecchio elettronico produce mediamente $ N $ impulsi al secondo. Sapendo che il numero di impulsi segue una legge di Poisson, stabilire la probabilità che in $ 1000 $ secondi si verifichino almeno $ 2500 $ impulsi;
\end{exe}

\section{Variabili aleatorie continue}

\begin{exe}
	Sia $ X $ una variabile aleatoria continua di densità:
	\[ F_{X}(t) = ct(1-t)^2 \mathds{1}_{(0,1)}(t).\]
	\begin{enumerate}
		\item Determinare il valore della costante $ c $ in modo tale che $ f_{X}(t) $ risulti una densità.
		\item Calcolare il valore atteso di $ X $.
		\item Calcolare l'espressione analitica della funzione di ripartizione $ F_{X}(t) $.
	\end{enumerate}
\end{exe}

\begin{exe}
	Calcolare il valore atteso della variabile aleatoria $ X $ avente come funzione di ripartizione:
	\[ F_{X}(t) = \begin{cases}
			1- \frac{1}{(1+t)^2} & \text{per} \ t>0   \\
			0                    & \text{per} \ t\le0 \\
		\end{cases} \]
\end{exe}

\begin{exe}
	Trovare il valore della costante $ k \in \R $ in modo che la funzione $ f : \R \to \R $ definita da:
	\[ f(x) := \begin{cases}
			0   & \mbox{per}\ x \le -2 \ \mbox{e}\ x \ge 1 \\
			k   & \mbox{per}\ -2<x<0                       \\
			1-x & \mbox{per}\ 0 \le x <1
		\end{cases} \]
	sia una densità di probabilità per la variabile aleatoria continua $ X $. Una volta trovato il valore di $ k $, determinare la funzione di distribuzione $ F(x) $ e calcolare la probabilità $ P(-1<X<2) $.
\end{exe}

\begin{exe}
	Trovare il valore del numero reale $ a > 0 $ tale che la funzione $ f : \R \to \R $ definita da
	\[ f(x) = \begin{cases}
			0   & \mbox{per} \ x < -1 \ \mbox{e} \ x \ge a, \\
			x^2 & \mbox{per} \ -1\le x \le 0,               \\
			1   & \mbox{per} \ 0 < x < a                    \\
		\end{cases} \]
	sia una densità di probabilità per la variabile aleatoria continua $ X $.
	\begin{enumerate}
		\item Una volta trovato il valore di a, determinare la funzione di distribuzione $ F(x) $;
		\item calcolare le probabilità $ P(X > 0)$ e $ P \left(-\frac{1}{2} \le X \le \frac{1}{2} \right)  $.
	\end{enumerate}
\end{exe}

\section{Distribuzione normale (o gaussiana)}

\begin{exe}
	Sia $ X $ una variabile aleatoria normalmente distribuita con media $ 3.5 $ e varianza $ 0.5625 $.
	Calcolare le probabilità:
	\begin{enumerate}
		\item $ P(X \le 5) $;
		\item $ P(X \ge 1.25) $.
	\end{enumerate}
\end{exe}
\begin{exe}
	Sia $ X $ una variabile aleatoria distribuita normalmente.
	\begin{enumerate}
		\item  Se $ X $ ha media $ \mu = 2.3 $ e deviazione standard $ \sigma = 0.08 $, calcolare la probabilità
		      \[ P(2.276 < X < 2.316). \]
		\item  Se ora X è un'altra variabile aleatoria distribuita normalmente tale che $ P (X < 2.5) = 0.5 $ e $ P(X < 2.55) = 0.6915 $, quanto valgono i valori $ \mu $ e $ \sigma $ relativi a $ X $?
	\end{enumerate}
\end{exe}
\begin{exe}
	L'altezza di un gruppo di $ 100 $ studenti universitari è distribuita normalmente con media
	$ \mu = 178\ cm $ e varianza $ \sigma^2 = 225\ cm^2 $. Qual è la probabilità che uno studente preso a caso fra questi $ 100 $ abbia statura.
	\begin{enumerate}
		\item  maggiore di $ 188\ cm $?
		\item  minore di $ 165 \ cm $?
		\item  compresa tra $ 173 $ e $ 183\ cm $?
	\end{enumerate}
\end{exe}
\begin{exe}
	Il punteggio ottenuto dagli studenti per una prova scritta di un esame universitario del
	primo anno di un corso di Scienze Naturali può essere modellizzato come una variabile aleatoria distribuita	normalmente con media $ 22 $ e varianza $ 9 $. Qual è la probabilità
	\begin{enumerate}
		\item  che uno studente ottenga un voto superiore al $ 25 $?
		\item che uno studente ottenga un voto compreso tra $ 21 $ e $ 24 $?
		\item che uno studente ottenga un voto inferiore al $ 18 $?
	\end{enumerate}
\end{exe}


\end{document}
